\subsubsection{2.3. Đặc trưng MFCC của âm thanh}
\indent\indent Trong xử lý tín hiệu âm thanh và nhận dạng tiếng nói, việc trích xuất đặc trưng là bước chuyển đổi tín hiệu miền thời gian sang miền tần số rút gọn, giúp mô hình loại bỏ các nhiễu nền và tập trung vào các thành phần mang thông tin ngữ nghĩa. Phương pháp phổ biến và hiệu quả nhất hiện nay là sử dụng Mel-Frequency Cepstral Coefficients (MFCC) kết hợp với các đặc trưng động của chúng.\vspace{0.3cm}

MFCC \cite{davis1980comparison} được thiết kế dựa trên đặc điểm sinh lý học của tai người. Tai người không cảm nhận tần số theo tuyến tính mà theo thang đo logarit (nhạy cảm hơn với sự thay đổi ở tần số thấp và kém nhạy cảm hơn ở tần số cao).\vspace{0.3cm}

Quy trình trích xuất MFCC bao gồm các bước: Tín hiệu âm thanh được chia thành các khung (frames) ngắn (thường 20-30ms), rồi áp dụng biến đổi Fourier, sau đó đi qua một bộ lọc (Filterbank) được phân bố theo thang đo Mel. Mối quan hệ giữa tần số Hertz ($f$) và tần số Mel ($m$) được mô tả qua công thức:\vspace{-0.1cm}
\[
    m = 2595 \log_{10}\left(1 + \frac{f}{700}\right)
\]

Cuối cùng, phép biến đổi cosine rời rạc được áp dụng lên log năng lượng của các bộ lọc để tạo ra các hệ số MFCC, giúp loại bỏ sự tương quan giữa các chiều dữ liệu.\vspace{0.3cm}

Điểm đặc trưng của MFCC là chỉ biểu diễn được các đặc trưng quang phổ tĩnh (static spectral features) của một khung tại một thời điểm, bỏ qua tính chất thay đổi theo thời gian của tín hiệu âm thanh. Để khắc phục điều này, Furui \cite{furui1986speaker} đã giới thiệu việc sử dụng các đạo hàm bậc một và bậc hai của chuỗi MFCC:\vspace{-0.1cm}
\begin{itemize}[label={--}, itemsep=1pt]
    \item Delta: Là đạo hàm bậc một của MFCC, biểu diễn vận tốc thay đổi của các hệ số cepstral giữa các khung liên tiếp. Công thức tính hệ số Delta $d_t$ tại thời điểm $t$ thường được tính bằng hồi quy tuyến tính trên một cửa sổ ngắn $N$:
    \[
        d_t = \frac{\sum_{n=1}^{N} n (c_{t+n} - c_{t-n})}{2 \sum_{n=1}^{N} n^2}
    \]
    Trong đó $c_t$ là hệ số MFCC tại thời điểm $t$.
    
    \item Delta-Delta: Là đạo hàm bậc hai của MFCC (hay đạo hàm của Delta), đại diện cho gia tốc hay sự thay đổi của Delta theo thời gian.
\end{itemize}

Trong thực nghiệm, một vector đặc trưng âm thanh hoàn chỉnh thường là sự kết hợp của ba thành phần này. Ví dụ, nếu sử dụng 13 hệ số MFCC đầu tiên, kết hợp với 13 hệ số Delta và 13 hệ số Delta-Delta, ta sẽ thu được một vector đặc trưng 39 chiều cho mỗi khung thời gian, nắm bắt đầy đủ cả thông tin phổ âm và động lực học của tín hiệu.